\documentclass[a4paper, 12pt]{article}
\usepackage[T2A]{fontenc}
\usepackage[english,russian]{babel}
\usepackage{amsmath}
\usepackage{amsthm}
\usepackage{amsfonts}
\usepackage{amssymb}
\usepackage{pgfplots}

\usepackage{enumitem}
% Абзацы и геометрия страницы
\usepackage{indentfirst}
%\usepackage{geometry}
\usepackage[
    top=2cm,
    bottom=2cm,
    left=1cm,
    right=1cm
]{geometry}
\begin{document}


\begin{center}
    МИНИСТЕРСТВО НАУКИ И ВЫСШЕГО ОБРАЗОВАНИЯ РОССИЙСКОЙ ФЕДЕРАЦИИ

    ГОСУДАРСТВЕННОЕ ОБРАЗОВАТЕЛЬНОЕ УЧРЕЖДЕНИЕ ОБРАЗОВАНИЯ

    РЯЗАНСКИЙ ГОСУДАРСТВЕННЫЙ РАДИОТЕХНИЧЕСКИЙ УНИВЕРСИТЕТ 

    ИМ. В. Ф. УТКИНА

    КАФЕДРА ВЫЧИСЛИТЕЛЬНОЙ И ПРИКЛАДНОЙ МАТЕМАТИКИ 
\end{center}

\vspace{5 cm}

\begin{center}
    ЛАБОРАТОРНАЯ РАБОТА № 3

    ПО ДИСЦИПЛИНЕ "ТЕОРЕТИЧЕСКАЯ ИНФОРМАТИКА"

    НА ТЕМУ "ПРЕДСТАВЛЕНИЕ ЦЕЛЫХ ЧИСЕЛ В ЭВМ. ЧАСТЬ 1"

    ВАРИАНТ 14

\end{center}

\vspace{7 cm}

\begin{flushright}
    Выполнил:

    Студент группы 6413

    Косяков Д.А.

    Проверил

    Доцент кафедры ВПМ

    Никичкин Борис Владимирович
\end{flushright}

\begin{center}
    РЯЗАНЬ 2026
\end{center}

\pagebreak

\textbf{ Цель работы:} посчитать сумму целых чисел представленных в двоичной системе.

\textbf{ Задание:} перевести числа А и В в десятичную систему счисления, 
сложить их и результат перевести напрямую в шестнадцатеричную систему счисления, минуя двоичную систему счисления.
числа А и В представлены в дополнительном коде.

\textbf{ Выполнение задания:} заданы $A_2$ = 00101110 и $B_2$ = 10010110. 

Переведем число А из двоичной системы счисления в десятичную. Так как знаковый бит равен 0 --- число положительно:
 $$A_{10} = 0 \cdot 2^7 + 0 \cdot 2^6 + 1 \cdot 2^5 + 0 \cdot 2^4 + 1 \cdot 2^3 + 1 \cdot 2^2 + 1 \cdot 2^1 + 0 \cdot 2^0 = 32 + 8 + 4 + 2 = 46$$

Переведем число В из двоичной систмы счисления в десятичную. Так как знаковый бит равен 1 --- число отрицательно, чтобы найти модуль инвертируем все биты и прибавляем 1. 
Инверсия $B_2$ = 01101001, инверсия $B_2$ + 1 = 01101010:
$$|B_{10}| = 64 + 32 + 8 + 2 = 106$$
$$B_{10} = -106$$

Вычисляем сумму чисел $A_{10}$ и $B_{10}$:
$$A_{10} + B_{10} = 46 + (-106) = -60$$
Переведем число $A_{10} + B_{10}$ в шестнадцатеричную систему:

$$60 : 16 = 3, \text{ остаток 12 } (C)$$
$$3 : 16 = 0,  \text{остаток } 3$$

Если $60_{10} = 3C_{16}$, значит $-60_{10} = -3C_{16}$.
\vspace{1cm}

\textbf{Отчет:}
\begin{enumerate}[nosep]
    \item А = 00101110, В = 10010110.
    \item $A_{10} = 46$, $B_{10} = -106$.
    \item $A_{10} + B_{10} = -60$.
    \item $-60_{10} = -3C_{16}$.
\end{enumerate}

\end{document}







